% 太阳系（综述）
% license CCBYSA3
% type Wiki

本文根据 CC-BY-SA 协议转载翻译自维基百科\href{https://en.wikipedia.org/wiki/Solar_System}{相关文章}。


太阳系（The Solar System）由太阳及其所环绕的天体组成。该名称来源于“*Sōl*”，即太阳的拉丁名称。大约在 46 亿年前，一个分子云中密度较高的区域发生坍缩，形成了太阳以及一个原行星盘（protoplanetary disc），环绕天体便由此汇聚而成。太阳核心内部的氢融合为氦，释放出能量，这些能量主要通过其外层光球（photosphere）向外辐射，从而在整个系统中形成一个由内而外逐渐降低的温度梯度。太阳系中超过 **99.86%** 的质量集中在太阳内部。

在绕太阳运行的天体中，质量最大的为八大行星。按与太阳距离递增的顺序，最靠近太阳的四颗类地行星（terrestrial planets）分别是——**水星（Mercury）**、**金星（Venus）**、**地球（Earth）** 和 **火星（Mars）**。它们构成了太阳系的**内行星系统（inner Solar System）**。其中，地球与火星是太阳系中仅有的两颗位于太阳**宜居带（habitable zone）**内的行星——在此范围内，行星表面可存在液态水。

在距离太阳约 **5 个天文单位（AU）** 之外、越过“霜线”（frost line）的位置，分布着四颗更为庞大的行星：两颗**气态巨行星（gas giants）**——**木星（Jupiter）**与**土星（Saturn）**，以及两颗**冰巨行星（ice giants）**——**天王星（Uranus）**与**海王星（Neptune）**。它们共同构成了太阳系的**外行星系统（outer Solar System）**。其中，木星与土星占据了太阳系除恒星外几乎 **90%** 的全部质量。
```
